% =====================================================================
% COSC 3337 — Homework 3 — ANSWER TEMPLATE for Part 1
%
% Typesetting Part 1 in LaTeX earns 5 bonus points.
%
% To use this:  upload it to Overleaf (or run pdflatex locally), put your name
% in the \student command, fill in the answers, and submit the PDF.
% You do not submit this .tex file.
%
% ---------------------------------------------------------------------
% WRITING CODE AND SQL
%
% Put code (and anything else you want shown exactly, character for character)
% in a code block:
%
%   \begin{lstlisting}
%   SELECT name FROM artist WHERE country = 'Spain'
%   \end{lstlisting}
%
% Inside a code block, everything is printed exactly as you type it: no
% escaping needed.
%
% For a short piece of code inside a sentence, use \lstinline|...| :
%
%   The query \lstinline|SELECT * FROM artist| returns every artist.
%
% Outside code, underscores must be written \_ , as in artist\_id.
% =====================================================================
\documentclass[11pt]{article}

\usepackage[margin=1in]{geometry}
\usepackage{amsmath,amssymb}
\usepackage{booktabs}
\usepackage{enumitem}
\usepackage{fancyhdr}
\usepackage{textcomp}
\usepackage{listings}

\lstset{
  basicstyle=\ttfamily\small,
  columns=fullflexible,
  keepspaces=true,
  upquote=true,
  showstringspaces=false,
  breaklines=true,
  frame=leftline,
  xleftmargin=1.2em,
}

\newcommand{\student}{YOUR NAME HERE}

\pagestyle{fancy}
\fancyhf{}
\lhead{\small COSC 3337 — HW 3 — \student}
\rhead{\small Page \thepage}
\renewcommand{\headrulewidth}{0.4pt}

\setlength{\parindent}{0pt}
\setlength{\parskip}{6pt}

% =====================================================================
\begin{document}

\begin{center}
  {\Large\bfseries COSC 3337 — Homework 3}\\[2pt]
  {\large \student}
\end{center}

% =====================================================================
\section*{1.1 Application architecture}

\textbf{(a)}
\begin{enumerate}
  \item ...
\end{enumerate}

\begin{lstlisting}
your SQL here
\end{lstlisting}

\textbf{(b)}

\textbf{(c)}
\begin{enumerate}
  \item ...
  \item ...
\end{enumerate}

% =====================================================================
\section*{1.2 User input and SQL}

\textbf{(a)}
\begin{lstlisting}
your SQL here
\end{lstlisting}

\textbf{(b)}
\begin{lstlisting}
your SQL here
\end{lstlisting}
Rows returned:

\textbf{(c)} Input:
\begin{lstlisting}
your input here
\end{lstlisting}
SQL MySQL receives:
\begin{lstlisting}
your SQL here
\end{lstlisting}

\textbf{(d)}

% =====================================================================
\section*{1.3 Queries in a loop}

\textbf{(a)}

\textbf{(b)}
\begin{lstlisting}
SELECT ...
\end{lstlisting}

\textbf{(c)}

% =====================================================================
\section*{1.4 Scaling out: sharding}

\textbf{(a)}
\begin{enumerate}
  \item ...
  \item ...
  \item ...
\end{enumerate}

\textbf{(b)}

% =====================================================================
\section*{1.5 MapReduce by hand}

\textbf{(a)} Map output:

\textbf{(b)} Shuffle output:

\textbf{(c)} Reduce output:

\textbf{(d)}
\begin{lstlisting}
SELECT ...
\end{lstlisting}

\textbf{(e)}

\end{document}
